\documentclass[11pt]{article}
\usepackage[T1]{fontenc}
\usepackage{lmodern}
\usepackage[margin=1in]{geometry}
\usepackage{amsmath,amssymb,amsthm,mathtools}
\usepackage{microtype}
\usepackage{booktabs,longtable}
\usepackage[hidelinks]{hyperref}
\newtheorem{theorem}{Theorem}[section]
\newtheorem{lemma}[theorem]{Lemma}
\newtheorem{proposition}[theorem]{Proposition}
\newtheorem{corollary}[theorem]{Corollary}
\newtheorem{claim}[theorem]{Claim}
\theoremstyle{definition}
\newtheorem{definition}[theorem]{Definition}
\newtheorem{example}[theorem]{Example}
\theoremstyle{remark}
\newtheorem{remark}[theorem]{Remark}
\DeclareMathOperator{\lcm}{lcm}
\DeclareMathOperator{\id}{id}
\title{A polynomial error estimate for superpermutation length}
\author{}
\date{26 October 2026}
\begin{document}
\maketitle
\begin{abstract}
A superpermutation on $n$ letters is a word containing every permutation
of those letters as a contiguous substring. Let $L(n)$ be its least possible
length. We prove
\[
 L(n)=n!+(n-1)!+(n-2)!+2(n-4)!
       +O\bigl((n-4)!n^{-1/61}\bigr).
\]
The quantitative step uses a finite collection of joint label and cut
states and a rational transfer argument. The latter gives the explicit
coefficient $2+74/r$ once $n\ge 8\cdot10^{12}\cdot352^4r^{61}$, for
$r\ge6$ divisible by six. The accompanying Lean 4.35.0 formalization proves
the finite covering-word bound and the two-sided asymptotic statement,
including the construction of the source families and their assembly.
\end{abstract}

\section{Words and the quantitative target}

A \emph{word} is a finite sequence of letters. A word on an alphabet of
size $n$ is a \emph{superpermutation} if every permutation of the alphabet
occurs in consecutive positions. For example, $121$ is a superpermutation
on two letters. Concatenating all permutations shows that a superpermutation
exists for every positive $n$, so its minimum length $L(n)$ is defined.
Put
\[
 F_3(n)=n!+(n-1)!+(n-2)!.
\]
All factorials in this manuscript have nonnegative arguments.

The main finite bound uses the absolute integer
$C=8\cdot10^{12}\cdot352^4$.

\begin{theorem}\label{thm-main-finite}
For every integer $r\ge6$ divisible by six and every integer
$n\ge Cr^{61}$, there is a superpermutation $w$ on $n$ letters such that
\begin{equation}\label{eq-main-finite}
 r|w|\le rF_3(n)+(2r+74)(n-4)!.
\end{equation}
\end{theorem}

We prove the theorem by constructing finite families of compatible rows
and then assembling them into a word. The source families, their unused
cyclic classes, and the auxiliary circles for assembly are all part of
the construction. The final estimate is obtained by normalizing their
exact integer counts.

\begin{corollary}\label{cor-main-asymptotic}
As $n$ tends to infinity through the positive integers,
\[
 L(n)=F_3(n)+2(n-4)!+
 O\!\left((n-4)!n^{-1/61}\right).
\]
\end{corollary}

The matching lower estimate supplied by the formal development of
\cite{halfar-asym} is
\begin{equation}\label{eq-main-lower}
 L(n)\ge F_3(n)+2(n-4)!-2(n-5)!\qquad(n\ge7).
\end{equation}
The precise imported declarations are \texttt{ConstAsym.lower2\_L} and
\texttt{ConstAsym.lower2\_ratio}. The complete dependency chain of these
results has been rebuilt under Lean 4.35.0.
We use \eqref{eq-main-lower} only for the final two-sided consequence.
Theorem~\ref{thm-main-finite} is established in
Corollary~\ref{cor-transfer-polynomial-threshold}. Selecting the largest
admissible multiple of six then proves
Corollary~\ref{cor-main-asymptotic}, as detailed in
Corollary~\ref{cor-transfer-asymptotic}.

We first describe the literal local moves and the restricted collection
of states that contains every permitted modification. Counting and
separating those states gives many mutually disjoint local constructions.
The next step combines constructions on overlapping sets of letters and
removes their possible collisions. The assembly section describes the
transport, the additional circles used for joining, and the extension by
an arbitrary number of letters. Finally, a rational estimate converts the
exact counts into \eqref{eq-main-finite}, and an integer choice of $r$
gives the exponent $1/61$.

\section{Literal contours and finite sets of states}
\label{geo:section}

This section constructs many families of cyclic bases that remain disjoint
after certain local changes. The local changes will later prevent
intersections between families built on different sets of control symbols.
The construction keeps track of the actual order of the symbols. This is
needed because the word assembly uses literal overlaps, not just cyclic
equivalence classes.

Fix an integer $r\geq 6$, and write $k=r-1$ and $d=r-3$. There are $M$
distinct \emph{background symbols}, in a fixed cyclic order, and $r$
distinct \emph{control symbols}. These two alphabets are disjoint. Number
the background symbols by the elements of $\mathbb Z/M\mathbb Z$, in their
cyclic order. A \emph{cut} is the position immediately before a background
symbol; its coordinate is the number of that symbol. There are $k$
stationary positions for control symbols and one mobile position. The
stationary positions have fixed identities $0,\ldots,k-1$. Their identities
do not change when their cuts pass the distinguished background symbol.

A \emph{literal base} is a list containing all $M+r$ symbols exactly once.
Two bases have the same \emph{cyclic class} if one is a cyclic rotation of
the other. We always write the mobile symbol last. After deleting it, the
remaining list is a representative of the stationary cyclic order. Thus
the stationary cuts and their resident control symbols, together with the
gap occupied by the mobile symbol, determine the literal base uniquely.
The gap is necessary: the stationary cuts and the assignment of control
symbols alone do not determine a base during a sequence of full moves.

Put $h=M+k$, so a base has length $h+1$. For a base
$b=b_0\cdots b_h$, define two literal operations by
\[
 F(b)=b_1\cdots b_{h-1}b_0b_h,
 \qquad
 S(b)=b_hb_0\cdots b_{h-3}b_{h-1}b_{h-2}.
\]
A \emph{full row} consists of a base with visible length $h+1$, and a
\emph{short row} consists of a base with visible length $h-1$. Their
\emph{charges} are respectively $0$ and $2$. A fixed extra symbol, called
the satellite, is outside the base alphabet and is attached to every row
when the rows are used in the word construction. It does not change in
this section.

The head of a row is the first $h-1$ symbols of its base. Its tail is the
first $h-1$ symbols after rotating the base left by one more than its
visible length. Two consecutive rows are compatible when the first tail
equals the second head. The displayed formulas show directly that the
successor of a full row with base $b$ may have base $F(b)$, and the
successor of a short row may have base $S(b)$. A \emph{closed contour} is a
nonempty cyclic list of compatible rows. Its source classes are the cyclic
classes of its bases, with one occurrence for each row.

\subsection{Positive updates and local simplicity}

Suppose a stationary position contains $A$, the background symbol at its
cut is $x$, and the mobile symbol is $T$. Apply $F$ until the base ends in
$A,x,T$, and then apply $S$. We call this a \emph{positive update} of that
stationary position. The cyclic order changes locally from
\[
 \cdots A,x,T\cdots
 \quad\hbox{to}\quad
 \cdots x,A,T\cdots.
\]
Consequently the old mobile symbol $T$ becomes the stationary resident at
the next background cut, and $A$ becomes mobile. The new mobile is
immediately before the new resident. In terms of the assignment of control
symbols to positions, the update exchanges the mobile position with the
chosen stationary position. No other stationary cut changes.

A \emph{batch} updates every stationary position once, in some prescribed
order. If the initial cuts are $q_0,\ldots,q_{k-1}$, then during batch $t$
the current cuts are
\begin{equation}
 q_i+t+\epsilon_i\pmod M,
 \qquad \epsilon_i\in\{0,1\}.
 \label{geo:cuts}
\end{equation}
Here $\epsilon_i$ is one exactly when position $i$ has already been
updated in the current batch. After the last update, increase $t$ by one
and reset all the $\epsilon_i$ to zero. The short row used for that last
update still has the old value of $t$ and has phase zero at its last
target. Only its successor uses the reset convention. Full moves change
the mobile gap and leave the data in \eqref{geo:cuts} unchanged.

The circular distance between two residues modulo $M$ is the smaller of
the two distances between them around the $M$ cuts. Call the initial cuts
\emph{robust} if the indexed family
\[
 0,\qquad q_i-q_j\quad(0\leq i,j<k,\ i\ne j)
\]
has circular distance at least five between any two different members.
The word ``indexed'' matters: repeated differences are forbidden. In
particular, the cuts are distinct. Perturbing each cut by zero or one
leaves distinct stationary cuts at circular distance at least four.

For example, with $r=6$ and $M=360$, the cuts
\[
 (q_0,q_1,q_2,q_3,q_4)=(0,5,25,125,155)
\]
are robust. Their positive integer differences are
$5,20,25,30,100,120,125,130,150,155$. These differences, their negatives
modulo $360$, and zero are separated by at least five.

\begin{lemma}
\label{geo:proper-run}
Assume the cuts are robust. After a positive update, the next positive
update reaches its required suffix in fewer than $h-1$ full moves.
\end{lemma}
\begin{proof}
The stationary list has $h$ distinct symbols. Each full move rotates it
one position, so the required target suffix is reached in a unique number
$u$ of moves with $0\leq u<h$. It remains to exclude $u=h-1$.

Let $x$ be the background symbol crossed by the preceding short move.
Immediately after that move, the stationary list ends in $x$, and the
position just updated has cut $x+1$. If $u=h-1$, its final symbol $x$
becomes its first symbol. On the other hand, let $z$ be the cut of the
next target. At the required suffix, the stationary list ends in that
target's resident followed by $z$. The successor cut $z+1$ is unoccupied,
because the current stationary cuts are separated by at least four.
Therefore the first stationary symbol at this suffix is the background
symbol $z+1$. We obtain $x=z+1$ modulo $M$.

The previous target now has cut $x+1=z+2$, whereas the next target has
cut $z$. If these are different stationary positions, their distance two
contradicts the separation just stated. If they are the same position,
the equality would give $2=0$ modulo $M$, which is impossible: robust
cuts include at least two different cuts at circular distance at least
five, so $M\geq10$. Thus $u\ne h-1$, proving the stronger bound.
\end{proof}

\begin{lemma}
\label{geo:simplicity}
Assume the cuts are robust. Consider a schedule of exactly $M$ batches
that returns the complete literal base to its initial value. Its source
classes are pairwise distinct.
\end{lemma}
\begin{proof}
Suppose two emitted rows have the same cyclic base. If their mobile
symbols are equal, they have $k$ common stationary control symbols;
otherwise they have $k-1$. In either case there are at least two. A common
stationary symbol occupies some position $i$ in the first row and some
position $j$ in the second. Equality of its background cut gives
\[
 t'-t=q_i-q_j+\epsilon_i-\epsilon'_j\pmod M.
\]
The ordered pairs $(i,j)$ for different common symbols are distinct. Two
off-diagonal pairs would place distinct nonzero ordered differences within
distance two. One diagonal and one off-diagonal pair would place a
nonzero difference within distance two of zero. Robustness excludes both
possibilities. Therefore every common symbol occupies the same stationary
position in both rows, and $t'-t$ is $0$, $1$, or $-1$ modulo $M$.

Choose the direction of comparison so that the rows lie in the same batch
or in consecutive batches of the cyclic schedule. Literal closure permits
this comparison across the last and first batches. In the same batch,
equality of the phases of a common stationary symbol implies that its
position is not updated between the rows. In consecutive batches, its
phase must be one in the first row and zero in the second. It has already
been updated before the first row, and is not updated again before the
second. Thus no common stationary position is updated in the interval
between the rows.

If the mobile symbols are equal, every stationary position is common.
The two rows lie in a single run of full moves, possibly including its
final short source row. The $h$ possible insertion gaps for the mobile
in the fixed stationary circle give $h$ different cyclic bases. By
Lemma~\ref{geo:proper-run}, that run has fewer than $h$ full moves, so no
two of its source rows have the same cyclic base.

If the mobile symbols differ, write them as $a$ and $b$. The only
stationary position not common to the two rows is one position $j$,
occupied by $b$ first and by $a$ second. All intervening updates must be
at $j$. There is at most one such update in one batch, and at most two
across the boundary between consecutive batches. The mobile changes, so
the number of these exchanges is odd. It is therefore exactly one.

Let $x$ be the cut of $j$ before that update. Before the short move, $b$
is stationary at $x$ and $a$ is mobile. Afterwards, $a$ is stationary at
$x+1$ and $b$ is mobile. Equality of the two cyclic bases requires $a$
to occupy cut $x+1$ before the update. There is only one such mobile gap,
because no stationary resident is at that cut before the update. This is
the gap of the short source row. The same equality requires the new
mobile $b$ to occupy cut $x$ afterwards. Immediately after the short move,
the stationary list begins with $a$ and ends with the background symbol
$x$. Inserting $b$ before $x$ therefore requires the final full gap,
reached after exactly $h-1$ full moves. Lemma~\ref{geo:proper-run} says
that the next target is reached strictly before this gap. The required
row therefore does not occur. This excludes the final possibility.
\end{proof}

\subsection{The normal order and exact replacement intervals}

Choose representatives $0=q_0<q_1<\cdots<q_{k-1}<M$. The \emph{normal
order} is
\[
 O=(k-1,k-2,\ldots,0).
\]
Begin with the mobile immediately before the resident at position zero.
As before, the mobile is the last symbol of the literal base. A complete
normal batch acts on the $r$ control positions as an $r$-cycle, denoted
$C$. Indeed, the original mobile goes to the first target; each subsequent
target receives the resident displaced from the preceding target; and the
resident of the last target becomes mobile.

\begin{lemma}
\label{geo:normal-counts}
If $r$ divides $M$, then $M$ normal batches form a literally closed
contour. Its number of rows $N_0$ and its charge $Q_0$ are
\[
 N_0=M\bigl((r-2)M+r(r-1)+1\bigr),
 \qquad Q_0=2(r-1)M.
\]
\end{lemma}
\begin{proof}
The role assignment returns because $C^M$ is the identity. Every
stationary cut advances $M$ times and therefore returns. The final target
is again zero, so the mobile is again immediately before its resident.
These data determine the same literal base, with the mobile last.

In one normal batch the first full run crosses $q_{k-1}+1$ background
symbols and all $k$ resident controls. After updating target $i$, the
full run to target $i-1$ crosses $M+q_{i-1}-q_i$ background symbols and
all $k$ resident controls. Consequently the total number of full moves is
\[
 q_{k-1}+1+k+
 \sum_{i=1}^{k-1}(M+q_{i-1}-q_i+k)
 =(k-1)M+1+k^2.
\]
There are also $k$ short rows, each of charge two. Multiplying the two
batch totals by $M$ gives the formulas.
\end{proof}

We next describe a replacement that keeps a specified control symbol out
of a specified stationary position for a chosen interval. Let $P$ be an
order of the $k$ stationary positions, with first target $a$. Performing
$P$ followed by its reverse is a \emph{parking pair}. Each position is
updated twice. The exchanges of roles cancel in reverse order, so every
control returns to its original role. The control that was initially
mobile is placed at $a$ by the first exchange and recovered by the last.
No intermediate exchange uses $a$. Throughout the pair, including all
full rows, this control is either mobile or resident at $a$.

For example, if $k=5$ and the forbidden position is $4$, use
$P=(3,4,2,1,0)$. A control initially mobile stays out of position $4$
throughout this order and its reverse. At the midpoint it is resident at
position $3$, although all stationary cuts have already advanced once.

\begin{lemma}
\label{geo:parking-endpoint}
Let $T$ be a positive integer. Replace $2rT+1$ consecutive normal batches
by $rT$ parking pairs, all with the same order $P$, followed by one normal
batch. The replacement has exactly the same final literal base as the
normal interval and the same total number of short rows. A control
initially mobile never occupies any stationary position different from
the first target of $P$ during the parking portion.
\end{lemma}
\begin{proof}
The parking pairs restore the role assignment, and the final normal
batch applies $C$. The normal interval also applies $C$, because
$C^{2rT+1}=C$. Both intervals advance every stationary cut $2rT+1$ times.
Both end with an update of position zero. The post-update mobile gap is
therefore the same, and the literal bases agree. Each interval has
$k(2rT+1)$ short updates. The final assertion follows from the preceding
description of each pair. Protection is asserted during the parking
portion; the final normal batch must lie outside the interval to be
protected.
\end{proof}

We call the final normal batch the \emph{cleanup batch}. It is needed
because a parking pair can end with its mobile immediately before a
resident other than that of position zero. Equality of role assignments
alone would not give equality of literal endpoints.

\subsection{A catalogue of joint role and phase states}

Only two parking orders are needed. Put $P_0=O$, and obtain $P_1$ by
exchanging the first two entries of $O$. For a forbidden position $j$,
use $P_1$ when $j$ is the first target of $P_0$, and use $P_0$ otherwise.
The chosen first target always differs from $j$. Fix this choice for all
parking pairs in a replacement interval.

An \emph{admissible schedule} is obtained from the normal $M$-batch
contour by replacing mutually disjoint intervals as in
Lemma~\ref{geo:parking-endpoint}, using only $P_0$ and $P_1$. An interval
may be placed on the cyclic schedule; choose the initial cut outside the
replacement intervals. Since each replacement has its normal literal
endpoint, all replacements can be performed without changing the states
at the other normal boundaries.

Let a role assignment be a bijection from the $r$ positions, including
the mobile position, to the $r$ named controls. Write $U_\ell(P)$ for the
role action of the first $\ell$ updates of $P$ followed by its reverse,
where $0\leq\ell\leq 2k$. For such a prefix define a phase vector by
\[
 \epsilon_\ell(P)=
 \begin{cases}
  \text{the indicator of the first $\ell$ targets of $P$},&0\leq\ell<k,\\
  \text{the indicator of the first $\ell-k$ targets of $P^{\rm rev}$},
       &k\leq\ell<2k,\\
  0,&\ell=2k.
 \end{cases}
\]
At the midpoint $\ell=k$, the translation has increased by one and the
phase vector is zero. The role action remains $U_k(P)$; it has not reset.
At $\ell=2k$, the translation has increased by two and both the phase
vector and the role action have reset.

Define the finite set of \emph{joint states}
\begin{equation}
 \mathcal J=
 \left\{\bigl(C^sU_\ell(P),\epsilon_\ell(P)\bigr):
      0\leq s<r,\ P\in\{P_0,P_1\},\ 0\leq\ell\leq2k\right\}.
 \label{geo:joint-set}
\end{equation}
Composition here means chronological action on a role assignment: the
normal batches act first, followed by the indicated prefix. We count
joint states, not independent choices of a role permutation and a phase
vector.

\begin{lemma}
\label{geo:joint-containment}
Fix an initial role assignment $\sigma$. Every source row of an admissible
schedule has stationary cuts $q_i+t+\epsilon_i$ and role assignment
$\sigma\tau$ for some translation $t$ and some
$(\tau,\epsilon)\in\mathcal J$. Moreover,
\[
 |\mathcal J|\leq2r(2k+1)=4r^2-2r<4r^2.
\]
\end{lemma}
\begin{proof}
A normal boundary has assignment $\sigma C^s$ for some $s$ modulo $r$.
The proper prefixes of its batch are first-half prefixes of the $P_0$
pair. Its final state is the empty prefix at normal phase $s+1$.

A replacement starts at a normal boundary. Every complete parking pair
restores its starting role assignment. Thus every pair in the replacement
has the same list of relative role actions $C^sU_\ell(P)$. The reset
convention in \eqref{geo:cuts} gives exactly the displayed phase vectors,
including the midpoint. The cleanup has the proper prefixes of a $P_0$
batch, and its final state is the next normal role phase. By
Lemma~\ref{geo:parking-endpoint}, the subsequent normal boundary agrees
with the original schedule, so the argument applies successively to all
replacements.

Full moves change neither the role assignment nor the stationary cuts.
The short row preceding the last update of a batch is represented by the
proper prefix before that update; its successor is represented at the
reset boundary. Thus every emitted row has been included. Finally, $r$
divides $M$, the translation increases by $M$, and the schedule closes
literally. The last and first states obey the same convention. The
cardinality bound counts the indices in \eqref{geo:joint-set}; repeated
representations only decrease the number of states.
\end{proof}

\subsection{Many robust modular frames}

The stationary cut set will have one additional algebraic property that
allows us to reconstruct a missing cut. Choose integers
\[
 0<z_1<\cdots<z_d<M,
 \qquad s=z_1+\cdots+z_d\pmod M,
\]
where $s$ is represented in $\{0,\ldots,M-1\}$. Sort the $k=d+2$ points
$0,z_1,\ldots,z_d,s$ into increasing order, and retain the result only
when it is robust. This ordered cut tuple is a \emph{modular frame}.
The point $s$ is its sum point. Assign coefficient $d-1$ to zero,
coefficient $1$ to the sum point, and coefficient $-1$ to every $z_i$.
These coefficients are all nonzero, their sum is zero, and their weighted
sum of cuts is zero modulo $M$. Thus the relation is unchanged by a
common translation. In the example above, $z=(5,25,125)$ and $s=155$.

We use the following elementary fact for arbitrary moduli. For a nonzero
integer $a$, the congruence $ax=b\pmod M$ has either no solutions or
$\gcd(|a|,M)$ solutions in $\mathbb Z/M\mathbb Z$. Indeed, if one solution
exists then all solutions form a translate of the kernel of multiplication
by $a$. On writing $g=\gcd(|a|,M)$ and dividing by $g$, coprimality shows
that the kernel consists of the $g$ multiples of $M/g$. In particular,
there are at most $|a|$ solutions. No primality assumption is needed.

\begin{lemma}
\label{geo:frame-supply}
Let $F_M$ be the number of robust modular frames. Then
\[
 F_M\geq
 \frac{(M-1)^d-18k^4(M-1)^{d-1}}{d!}
 \geq\frac{M^d}{d!}\left(1-\frac{19r^4}{M}\right).
\]
Each retained anchored cut set has a unique sum point.
\end{lemma}
\begin{proof}
First choose an ordered vector from $\{1,\ldots,M-1\}^d$, without sorting
its coordinates. Regard the cuts as evaluations of the formal points
$0,e_1,\ldots,e_d,\boldsymbol 1$ in $\mathbb Z^d$. Their affine relations
form a one-dimensional space, generated by
\[
 (d-1,-1,\ldots,-1,1).
\]
To verify this, a relation with coefficient $a$ at $\boldsymbol 1$ has
coefficient $-a$ at each $e_i$, by comparing coordinates, and coefficient
$(d-1)a$ at zero, by summing coefficients. Every nonzero relation
therefore involves all $d+2\geq5$ points.

A comparison of two distinct ordered differences, or of a nonzero
difference with zero, involves at most four formal points. Its linear
form is therefore nonzero. Every coefficient in that form has absolute
value at most two. There are at most $k^4$ comparisons. For robustness,
exclude the nine residues $-4,-3,\ldots,4$ for each comparison. Fix all
but one coordinate with nonzero coefficient. The modular root bound gives
at most two choices for the remaining coordinate for each excluded
residue. The number of bad ordered vectors is at most
$18k^4(M-1)^{d-1}$.

The good vectors are invariant under coordinate permutation, and a good
vector has distinct coordinates. Its orbit under coordinate permutation
therefore has exactly $d!$ elements. Sorting gives the first inequality.
For the second, use
$(M-1)^d\geq M^d-dM^{d-1}$ and
$d+18k^4\leq19r^4$.

It remains to check that sorting has not identified different sum
decorations. If an anchored cut set had two sum points $s$ and $s'$, the
sum of all its cuts would be both $2s$ and $2s'$ modulo $M$. Distinct
solutions require $M$ even and $s-s'=M/2$. The two ordered differences
$s-s'$ and $s'-s$ would then coincide. This contradicts robustness, so
the sum point is unique.
\end{proof}

\subsection{Reconstruction and disjoint labelled families}

A \emph{labelled frame} is a pair $(q,\sigma)$, where $q$ is a robust
modular frame and $\sigma$ is an initial bijective assignment of the $r$
named controls to the $r$ positions. Its \emph{envelope} is the set of
cyclic bases obtained by choosing an arbitrary translation, a joint state
$(\tau,\epsilon)\in\mathcal J$, and an arbitrary mobile insertion gap,
with stationary cuts $q_i+t+\epsilon_i$ and assignment $\sigma\tau$.
Lemma~\ref{geo:joint-containment} shows that every source row of every
admissible repair of this labelled frame lies in its envelope.

\begin{lemma}
\label{geo:reconstruction}
Fix a labelled frame and fix a joint state on each side of a possible
envelope intersection. The number of positional reconstruction witnesses
for a target frame is at most
\[
 D_0=k(k-1)^2\bigl(d+1+\gcd(d-1,M)\bigr)
       \leq2d\,k(k-1)^2\leq2r^4.
\]
A witness includes the matching of the retained stationary positions.
\end{lemma}
\begin{proof}
A common cyclic base has at least $k-1$ controls stationary on both sides.
If the mobile symbols agree, discard one of the $k$ common stationary
controls. Choose the omitted stationary position on the source side and
on the target side. The remaining $k-1$ stationary controls occur in the
same cyclic order. The binary phase perturbations do not change their
order, because the cuts remain separated by at least four. There are
therefore $k-1$ possible cyclic matchings.

For a matched source position $i$ and target position $j$, equality of
the cuts gives
\[
 q'_j+\delta=q_i+\epsilon_i-\epsilon'_j\pmod M,
\]
where $\delta$ is the difference of the two common translations. The
phase vectors are fixed, so the right-hand side is known. Choose also
the target sum position $a\in\{1,\ldots,k-1\}$. This fixes its coefficient
vector $c$. Set $w_j=q'_j+\delta$. Since $\sum c_j=0$, the target relation
is $\sum c_jw_j=0$. It determines the omitted coordinate $w_j$ in at most
$\gcd(|c_j|,M)$ ways. If position zero was retained, its known coordinate
$w_0$ determines $\delta$; if it was omitted, its reconstructed coordinate
does so. Subtracting $\delta$ then determines every target cut. Sorting
and robustness can reject a reconstruction but cannot add one.

For each fixed sum position, summing the root multiplicities over the
omitted target position gives
$d+1+\gcd(d-1,M)\leq2d$. There are $k$ source omissions, $k-1$ cyclic
matchings, and $k-1$ choices of sum position. Multiplying gives the stated
bound, which counts witnesses even if several reconstruct the same frame.
\end{proof}

\begin{lemma}
\label{geo:labelled-packing}
Put
\[
 \Delta=64r^8,
 \qquad \eta=\frac{r!}{2d!(\Delta+1)}.
\]
For every $M\geq38r^4$, there are
$\lfloor\eta M^d\rfloor$ labelled frames whose envelopes are pairwise
disjoint.
\end{lemma}
\begin{proof}
Form a finite graph whose vertices are all labelled frames, with an edge
when two distinct vertices have intersecting envelopes. Distinct vertices
with coincident envelopes are also adjacent. By
Lemma~\ref{geo:frame-supply}, the graph has at least
$r!M^d/(2d!)$ vertices.

Fix a source vertex. For a chosen pair of joint states, each witness of
Lemma~\ref{geo:reconstruction} determines the target frame and matches
$r-2$ named stationary control symbols. Since the target role action is
a permutation, those equalities determine the target initial assignment
on $r-2$ positions. The remaining two labels have at most two assignments
to the remaining positions. Thus the number of possible neighbouring
vertices, including the source vertex itself, is at most
\[
 2|\mathcal J|^2D_0\leq2(4r^2)^2(2r^4)=64r^8.
\]
This bound includes all assignments at the time the graph is formed. It
would not be valid to select unlabelled frames first and then attach all
$r!$ assignments to each selected frame.

Greedily select a vertex and remove it and all its neighbours. Each
selection removes at most $\Delta+1$ vertices. Continuing until no
vertices remain produces an independent set with at least
$r!M^d/(2d!(\Delta+1))=\eta M^d$ vertices. Select the required integer
number from this independent set. Their envelopes are disjoint by the
definition of the graph.
\end{proof}

For fixed $M$, the selected family can be used for every background circle
and every set of $r$ active letters. To see this, first make the selection
using $M$ formal background symbols and $r$ different formal control
symbols. Replace the background symbols, in order, by the letters of any
chosen background circle, and replace the formal controls by its active
letters. These substitutions give an injective map on the entire formal
alphabet. They commute with the construction of each row and with cyclic
rotation. If two substituted rows were equal up to rotation, injectivity
would give the same equality before substitution. Thus the selected
envelopes remain disjoint. In particular, the choice of frames is
independent of the guard positions introduced in the next section.

\subsection{What the local construction supplies to assembly}

For later use, the literal displacement also has an exact description.
Give each background symbol its index modulo $h$ in a literal base. A
full move changes that index by $-1$. A short move changes it by $+1$,
except for the crossed background symbol, which stays at index $h-1$.
In $M$ complete batches each stationary position crosses every background
symbol exactly once. Every background symbol therefore has exactly $k$
exceptional short moves. If the contour is literally closed, with $N$
rows and total charge $Q$, this proves
\begin{equation}
 Q-N\equiv k\pmod h,
 \qquad \gcd(|Q-N|,h)=\gcd(k,h)\leq k.
 \label{geo:displacement}
\end{equation}

An untouched normal batch contains an actual consecutive run of literal
$F$ successors of length at least
\[
 \frac{(k-1)M+k^2+1}{k}.
\]
It is an untouched batch that gives this estimate. The first full run of
a cleanup batch can start at a different mobile gap and is not covered
by the normal count. The global repair argument must therefore leave an
untouched batch before applying a word assembly lemma that requires such
a run.

The results in this section construct disjoint envelopes for one active
control alphabet on one fixed background. They also prove that every
admissible repair stays inside its envelope. To use several overlapping
active alphabets, one must additionally construct the forbidden-event
catalogue, prove the deletion estimates, and place disjoint replacement
intervals. An actual word bound further requires the lists of unused
source classes, an iteration that inserts a variable number of controls,
and additional cyclic words that meet every required literal overlap.
Those are separate mathematical statements; the numerical counts in this
section do not assume their existence.

\section{Combining the local constructions}
\label{sec-global}

Fix an integer $r\ge6$ divisible by six, and put $k=r-1$ and $d=r-3$.
We use the labelled frame family of the preceding section, with
\[
 \Delta=64r^8,\qquad
 \eta=\frac{r!}{2d!(\Delta+1)},\qquad
 \beta=(r-2)\eta.
\]
At each sufficiently large background size $M$, select exactly
$\lfloor\eta M^d\rfloor$ labelled frames whose allowed sets of cyclic
rows are pairwise disjoint. The selection is fixed before the modifications
described below. Define
\begin{equation}\label{eq-global-parameters}
 A=\left\lceil\frac{r^3}{\beta}\right\rceil,\quad
 Q_{
 \mathrm{grid}}=rA,\quad W=r^2A,\quad H=rA^2,\quad
 G=W-r,\quad \rho=\frac{A\beta}{r},\quad
 E=10^{12}r^{22}W^4.
\end{equation}
The letter $Q_{\mathrm{grid}}$ denotes a grid width. It is unrelated to
the row charge used in the assembly section.

\subsection{Sets of active letters and fixed background letters}

Take the $W$ control letters to be the pairs $(i,j)$ with
$0\le i<r$ and $j\in\mathbb Z/Q_{\mathrm{grid}}\mathbb Z$.
For $0\le a<A$ and $b\in\mathbb Z/Q_{\mathrm{grid}}\mathbb Z$, set
\[
 T_{a,b}=\{(i,ai+b):0\le i<r\}.
\]
The members of $T_{a,b}$ are the \emph{active letters} of that
construction. The other $G$ control letters are its \emph{guards}; these
letters stay in the background circle and are never chosen as mobile
letters. For example, when $r=3$ and $A=2$, the blocks $T_{0,0}$ and
$T_{1,0}$ in a grid with six columns share just $(0,0)$.

There are $H$ blocks, each with $r$ letters. Distinct blocks intersect
in at most one letter. Indeed, two common letters in distinct rows $i,j$
would imply
\[
 Q_{\mathrm{grid}}\mid(a-a')(i-j),\qquad
 |(a-a')(i-j)|\le(A-1)(r-1)<Q_{\mathrm{grid}}.
\]
The product must therefore be zero. Since $i\ne j$, this gives $a=a'$,
and either common letter then gives $b=b'$. This argument uses integers
and remains valid when $Q_{\mathrm{grid}}$ is composite. Each letter
belongs to $A$ blocks. It follows that a block meets exactly
$I=r(A-1)\le W/r$ other blocks: each of its letters lies in $A-1$ other
blocks, and no block is counted at two letters.

Fix a cyclic order of $m$ old letters. Insert each guard in a different
old cut, immediately before the old letter at that cut. A \emph{guard
background} is the resulting cyclic order of $M=m+G$ letters. There are
exactly
\[
 (m)_G=m(m-1)\cdots(m-G+1)
\]
guard backgrounds for each block. Distinct cuts give no ordering
multiplicity. For $m\ge G$,
\begin{equation}\label{eq-guard-count}
 1-\frac{G(G-1)}{2m}\le\frac{(m)_G}{m^G}\le1.
\end{equation}
The lower bound follows by applying
$\prod_i(1-u_i)\ge1-\sum_i u_i$ to $u_i=i/m$.

A \emph{candidate} consists of an active block, a guard background, and
one selected labelled frame. Candidates on the same block but different
backgrounds cannot share a row: deleting their active letters would recover
the same background from that row. Candidates on the same background are
already separated by the selected frame family. Thus only candidates on
different blocks can collide.

\subsection{A finite list of necessary collision equations}

Anchor every cyclic background at the same old letter. A local background
cut $y\in\{0,\ldots,M-1\}$ projects to an original cut
$x\in\{0,\ldots,m-1\}$. Counting the guards passed before $y$ gives
\[
 y\equiv x+e\pmod m,\qquad 0\le e\le G.
\]
This also covers the guard immediately before the anchor. In that case
the local cut has the form $m+e$ and the original cut is zero.

Every stationary frame satisfies a relation whose coefficient multiset is
\[
 (d-1,1,\underbrace{-1,\ldots,-1}_{d\text{ entries}}).
\]
The coefficients sum to zero, and their absolute values sum to $2d$.
Including the binary cut phases and reducing the local relation modulo
$m$ shows that any foreign row must satisfy
\begin{equation}\label{eq-foreign-relation}
 \sum_u c_u x_u\equiv b\pmod m,\qquad
 |b|\le C_1=(4d+1)G+2d\le4rW.
\end{equation}
For completeness, choose the canonical local cuts in $[0,M)$.
Their weighted sum differs from the binary phase sum by an integer
multiple of $M$, whose multiplier has absolute value at most $2d+1$.
Replacing $M$ by $m+G$ and each local cut by $x_u+e_u$ gives the bound
$(2d+1)G+2dG+2d$ displayed above.
There are $k(k-1)$ assignments of this coefficient multiset to the
stationary physical labels, since one chooses the label with coefficient
$d-1$ and then the label with coefficient $1$.

First delete a candidate if a relation \eqref{eq-foreign-relation} holds
for a possible foreign stationary set containing none of its active
letters. Such a relation has a nonzero coefficient on a guard variable.
There are at most
\[
 J_1\le Hr\,k(k-1)(2C_1+1)\le9rW^3
\]
relations to consider. A nonzero coefficient has absolute value at most
$r$. After all other guard variables are fixed, its congruence has at most
$r$ solutions modulo $m$, because the number of roots of $ax=b$ is at most
$\gcd(a,m)\le|a|$. Thus each equation excludes at most $rm^{G-1}$
guard tuples. If $m\ge G^2$, \eqref{eq-guard-count} gives
$(m)_G\ge m^G/2$. This first deletion removes at most the fraction
\begin{equation}\label{eq-guard-loss}
 \frac{18r^2W^3}{m}.
\end{equation}
The selected frame count cancels from this calculation.

Suppose two surviving candidates on different blocks share a cyclic row.
Their blocks must have a common letter $v$, and $v$ must be stationary in
both rows. If the blocks were disjoint, all foreign stationary letters
would be guards. If $v$ were mobile in one row, its stationary letters
would all be guards for the other candidate. Either case contradicts the
first deletion, applied to the appropriate candidate.

If $v$ occupies stationary slot $j$ during batch $t$, then its local cut
is $(q_j+t+\epsilon_j)\bmod M$. Here $0\le t<M$ is the canonical
batch number. The local wrap is zero or one, so its original cut obeys
\[
 x_v\equiv q_j+t+e'\pmod m,\qquad -2G\le e'\le1.
\]
Substituting in \eqref{eq-foreign-relation} gives the following finite
list of possible events:
\begin{equation}\label{eq-event}
 a(t+q_j)+\sum_u b_ux_u\equiv b_1\pmod m,
 \qquad |b_1|\le C_2=(6d-1)G+3d-1\le6rW.
\end{equation}
An \emph{event} forbids the pair $(v,j)$ for that entire batch.
The list ranges over every meeting foreign block, every foreign mobile
letter, every coefficient assignment, every candidate slot $j$, and every
integer $b_1$ in the stated interval. It therefore includes both the
original rows and every permitted later modification.

There are at most
\[
 J_{\mathrm{event}}\le I k^3(k-1)(2C_2+1)\le13r^4W^2
\]
equation types. For a fixed type there are at most $2|a|$ actual roots
$0\le t<M$, provided $M<2m$. For a fixed shared label, summing $|a|$
over its coefficient assignments gives $2d(d+1)$. Hence the total number
of events, with repetitions allowed, is at most
\begin{equation}\label{eq-event-count}
 K\le4Ik^2d(d+1)(2C_2+1)\le52r^4W^2.
\end{equation}

\subsection{Separating events without a distribution assumption}

Put $L=2r+1$ and $S=2L$. Delete a candidate if it has two events for
different forbidden pairs at circular distance at most $S$ on the
$M$-batch circle. Eliminating the two event times gives
\begin{equation}\label{eq-proximity}
 aa'(q_j-q_l)+\sum_u(a'b_u-ab'_u)x_u
 \equiv a'b_1-ab'_1+aa'(\delta+\sigma G)\pmod m,
\end{equation}
where $|\delta|\le S$ and $\sigma\in\{-1,0,1\}$.
The term $\sigma G$ accounts for a wrap on the $M$-circle. It must be
retained even though the congruence is modulo $m$.
There are at most
\[
 J_2\le3(2S+1)J_{\mathrm{event}}^2\le4563r^9W^4
\]
such equations.

Every equation \eqref{eq-proximity} has a nonzero variable coefficient.
Here are all possibilities. If the foreign blocks differ, each has
$r-2$ guard letters in its stationary set, and the two blocks share at
most one letter. A private guard then has a nonzero coefficient. If the
foreign block is the same but its mobile letters differ, a guard present
in only one stationary set has a nonzero coefficient. If the block and
mobile agree, either a guard coefficient $a'b_u-ab'_u$ is nonzero or the
two full coefficient vectors are proportional. In the latter case their
equal multisets force the ratio to be $1$. Indeed their squared norms
force the ratio to be $1$ or $-1$, and $-(d-1)$ does not occur in the
multiset when $d\ge3$. Thus the vectors are equal. The forbidden pairs
then differ only in their slots, so the remaining form is
$a^2(q_j-q_l)$ with $j\ne l$. On a fixed position of the sum point,
the cuts have the distinct coefficient vectors
$0,e_1,\ldots,e_d,\boldsymbol{1}$. Their difference is nonzero.
All coefficients have absolute value at most $3r^2$.

We count the guard-variable and guard-free equations separately. If a
guard coefficient is nonzero, fix the selected labelled frame and all
other guards. There are at most $3r^2$ roots modulo $m$. Dividing by
$(m)_G\ge m^G/2$ bounds the deleted fraction by $6r^2/m$ per equation,
uniformly in the selected frame.

If all guard coefficients vanish, partition the sorted frame parameters
according to the position of their sum point and the integer
\[
 \ell=\left\lfloor\frac{\sum_i z_i}{M}\right\rfloor.
\]
There are at most $d(k-1)\le r^2$ choices. Within such a part,
$q_j=v_j\cdot z-\mathbf{1}_{j=a}\ell M$. The winding term changes only
the constant in \eqref{eq-proximity}. Fix a coordinate with nonzero
coefficient. It ranges in an interval shorter than $2m$, so it has at
most $6r^2$ roots. Deleting that coordinate leaves a strictly increasing
list of $d-1$ integers, with at most $M^{d-1}/(d-1)!$ possibilities.
Including all $r!$ initial assignments and dividing by
$\lfloor\eta M^d\rfloor\ge\eta M^d/2$ bounds the deleted fraction,
per part and per equation, by
\[
 \frac{12r^2r!}{\eta(d-1)!M}
 =\frac{24r^2d(\Delta+1)}{M}.
\]
This counts in the entire parameter set before restricting to the selected
frames. It requires no uniform distribution property of that selection.
Since $\Delta+1\le65r^8$ and $M\ge m$, both cases together are bounded by
\begin{equation}\label{eq-proximity-loss}
 \frac{7\,118\,280r^{22}W^4}{m}.
\end{equation}

\subsection{Modifying every surviving candidate}

Around every event take the closed interval of radius $L$ on the batch
circle, and merge intersecting intervals. Each connected component has
one forbidden pair, since intervals for different pairs cannot intersect.
If there are no events, keep the normal construction. Otherwise,
$M>(2L+1)K$ ensures a gap longer than $2L+1$ between consecutive event
times. Cut in this gap and regard the remaining components as ordinary
intervals. The cut changes only where these intervals are written; the
canonical times in \eqref{eq-event} remain unchanged.

For a component whose first and last events are $f,e$, let $(v,j)$ be
its forbidden pair. Choose the normal boundary $s\in[f-r+1,f]$ at
which $v$ is mobile; the normal role action is a cycle on the $r$ roles.
Set
\[
 P=2r\left\lceil\frac{e+1-s}{2r}\right\rceil.
\]
Replace the next $P$ normal batches by $P/2$ pairs of a permitted order
and its reverse, using the order whose first target is different from
$j$. Follow them by one normal cleanup batch. The local endpoint theorem
shows that these $P+1$ batches end at exactly the normal literal state.
Throughout the first $P$ batches, the letter $v$ is either mobile or in
the first target slot. Thus it never occupies $j$ during any event batch.

The inequalities
\[
 f-r+1\le s\le f,\qquad
 e+1\le s+P\le e+2r,\qquad s+P+1\le e+2r+1
\]
put the entire replacement inside its merged interval. Different
replacements are disjoint. Each restores its normal endpoint, so they can
all be performed on the same cyclic walk. The joint-state containment
proved earlier keeps every modified row inside the originally selected
allowed set. If two modified candidates met, their unique shared letter
would be stationary in both rows and would produce an event
\eqref{eq-event}. The modification prevents that occupancy. This excludes
collisions between two modifications as well as collisions with normal
rows.

Regard a batch starting at integer time $t$ as the unit interval $[t,t+1]$.
The inequalities above put every changed batch inside the interval union.
That union has total length at most $2LK\le5rK$, so at most $5rK$
batches change. A complete batch has
at most $k(M+k)$ source rows. Comparing each changed batch with its normal
version therefore changes the total row count by at most $10r^2KM$ when
$M\ge r$. The short charge stays exactly $2(r-1)M$. There is also an
untouched normal batch. Its full rows follow the literal successor $F$,
which is the additional hypothesis needed in the connector construction.

\subsection{Uniform counts}

Let $B$ be the total number of source rows in the surviving modified
contours, let $q$ be their total short charge, and let $c_{\mathrm{adv}}$
be the number of circles in the actual connector lists constructed in the
assembly section. Put
$D_{m,W}=\prod_{i=0}^{W-1}(m+i)$ and $c_r=2(r-1)/(r-2)$.

\begin{proposition}\label{prop-global-counts}
For every multiple $m$ of $r$ with $m\ge4E$, the construction has literal
closure, pairwise distinct source cyclic classes above the chosen old
cyclic order, and the safe connector lists of the assembly section. Its counts satisfy
\[
 \frac{B}{D_{m,W}}=\frac{\rho W}{m}(1+\theta),\qquad
 \frac{q+c_{\mathrm{adv}}}{D_{m,W}}
   =\frac{c_r\rho W}{m^2}(1+\epsilon),\qquad
 |\theta|,|\epsilon|\le\frac Em,
\]
and $c_{\mathrm{adv}}\le(r-1)q/m$.
\end{proposition}

\begin{proof}
We first verify the size conditions used in the construction. The ceiling
defining $A$ gives
\[
 \eta^{-1}=\frac{r-2}{\beta}\le\frac{r}{\beta}
 \le\frac{W}{r^4}\le W.
\]
Also $W\ge r^2\ge1$ and $G\le W$. The definition of $E$ implies
$E\ge W^2$, $E\ge rW$, and $4E\ge38r^4$. Hence
\[
 M\ge m\ge38r^4,\qquad
 m\ge4W^2\ge G^2,\qquad
 G<m,\qquad M=m+G<2m.
\]
Since $d\ge3$ and $M\ge m\ge4W$, we have
$\eta M^d\ge M/W\ge4$. Moreover,
$rG/m\le rW/(4E)\le1/4$. Finally, the event bound gives
\[
 (4r+3)K\le260r^5W^2<4\cdot10^{12}r^{22}W^4\le m\le M.
\]
The strict middle inequality follows from $r,W\ge1$ and
$260<4\cdot10^{12}$. These prove all the required size conditions.
Also $r\mid M$, since $r$ divides $m$ and $W$. The local closure and
distinctness results and the preceding modifications consequently apply.
Every source class is above the chosen old cyclic order. Deleting all
controls first removes the guards from the guard background. A full move
only rotates the remaining old cyclic order. A positive short update
moves a control past one background letter, so it does not change the
relative cyclic order of the old letters. This remains true during every
replacement. Deleting all controls from any resulting row therefore
recovers the chosen old class, as required by
Lemma~\ref{lem-assembly-universe}.

Each modified contour has
\[
 N=(r-2)M^2+[r(r-1)+1]M+\delta_N,
 \quad |\delta_N|\le10r^2KM,\quad q_0=2(r-1)M.
\]
Let $h=M+r-1$ be the source port modulus and let $g$ be its greatest
common divisor with the monodromy. The local result gives $1\le g\le r-1$.
Write $k=r-1$. The untouched normal batch has $k$ short rows and
$(r-2)M+r(r-1)+1$ rows in total. One of its $k$ consecutive full runs
therefore has length
\[
 R\ge\frac{(k-1)M+k^2+1}{k}.
\]
Its bases follow the literal successor $F$, and the local construction
gives $R\le h$. To check the numerical hypothesis of the connector
lemma, use $g\lfloor R/g\rfloor\ge R-g$ and $g\le k$. These give
\[
 2g\lfloor R/g\rfloor
 \ge2R-2g
 \ge\frac{2(k-1)M+2}{k}
 \ge M+k=h.
\]
For the last inequality, it suffices that $(k-2)M+2\ge k^2$.
This follows from $M\ge2k$ and $k\ge5$, since
$2k(k-2)+2\ge k^2$. Thus the connector lemma supplies an actual list
of at most $2g^2\le2(r-1)^2$ circles for each contour.
Summing over the contours gives
$c_{\mathrm{adv}}\le(r-1)q/M\le(r-1)q/m$.

Before deletion the number of candidates is
$C_0=H(m)_G\lfloor\eta M^d\rfloor$. Define
\[
 \phi=\frac{(m)_G}{m^G},\qquad
 \psi=\frac{\lfloor\eta M^d\rfloor}{\eta M^d},\qquad
 \mu=\frac{M}{m},\qquad
 \chi=\frac{D_{m,W}}{m^W}.
\]
The exact leading contributions before deletion are
\begin{align*}
 C_0(r-2)M^2
   &=H\beta m^{W-1}\phi\psi\mu^{d+2},\\
 C_0q_0
   &=c_rH\beta m^{W-2}\phi\psi\mu^{d+1}.
\end{align*}
The exponents follow from $G+d+2=W-1$ and $G+d+1=W-2$.
Let $t$ be the proportion of candidates that survive both deletions.
The deletion estimates give
\[
 1-\frac{18r^2W^3+7\,118\,280r^{22}W^4}{m}\le t\le1.
\]
In particular $t>0$ under the present size bound. Averaging the relative
row errors over the surviving contours gives a number $u$ such that
\[
 |u|\le\frac{2r+1040r^5W^2}{m}.
\]
Their short charges are all exactly $q_0$. Write
$v=c_{\mathrm{adv}}/q$, so $0\le v\le(r-1)/m$. The two normalized
counts now have the exact expressions
\begin{align*}
 \frac{B}{D_{m,W}}
   &=\frac{H\beta}{m}
     \frac{\phi\psi t\mu^{d+2}(1+u)}{\chi},\\
 \frac{q+c_{\mathrm{adv}}}{D_{m,W}}
   &=\frac{c_rH\beta}{m^2}
     \frac{\phi\psi t\mu^{d+1}(1+v)}{\chi}.
\end{align*}
The relative-error numerators, each to be divided by $m$, are bounded by
\[
\begin{array}{lr}
\text{rounding the selected frame count}&W\,,\\
\text{distinct guard positions}&W^2/2\,,\\
\text{guard-only deletion}&18r^2W^3\,,\\
\text{nearby incompatible events}&7\,118\,280r^{22}W^4\,,\\
\text{row count within a contour}&2r+1040r^5W^2\,,\\
\text{additional connector charge}&r\,,\\
\text{replacing powers of }M\text{ by powers of }m&2W^2\,,\\
\text{the product }D_{m,W}/m^W&W^2\,.
\end{array}
\]
The first bound uses $d\ge3$ and $\eta^{-1}\le W/r^4$.
The next three were proved in \eqref{eq-guard-count},
\eqref{eq-guard-loss}, and \eqref{eq-proximity-loss}.
The row bound follows by dividing the displayed estimate for $N$ by
$(r-2)M^2$. We use the product inequalities
\[
 \prod_i(1+u_i)\le\frac1{1-\sum_i u_i},\qquad
 \prod_i(1-u_i)\ge1-\sum_i u_i
 \quad(0\le u_i,\ \sum_i u_i<1).
\]
For the first inequality, use $1+u_i\le1/(1-u_i)$ and the second
inequality. For the powers of $\mu$, there are at most $r$ equal factors
$1+G/m$. Their deviations sum to at most $rG/m\le1/4$, so
\[
 1\le\mu^{d+1}\le\mu^{d+2}
 \le\frac1{1-rG/m}\le1+\frac{2rG}{m}
 \le1+\frac{2W^2}{m}.
\]
For $\chi$, the deviations $i/m$ sum to at most $W^2/(2m)\le1/8$,
which gives $1\le\chi\le1+W^2/m$ in the same way.

The sum of the listed numerators is less than $9\cdot10^6r^{22}W^4$.
Let $s$ be that sum divided by $m$. Then $s<1/1000$. To combine the
estimates, let $s_0$ be the sum of the deviations allowed for the
numerator factors in either exact expression, and let $s_1$ be the
deviation allowed for $\chi$. Thus $s_0+s_1\le s$.
The product inequalities give a numerator between $1-s_0$ and
$1/(1-s_0)$, while $1\le\chi\le1+s_1$. The ratio is therefore at
least
\[
 \frac{1-s_0}{1+s_1}\ge1-s_0-s_1\ge1-s
\]
and at most $1/(1-s_0)\le1/(1-s)\le1+2s$.
In particular its deviation from one is at most $4s$.
Since $36\cdot10^6<10^{12}$, each relative error is at most $E/m$.
Finally $H\beta=\rho W$, as required.
\end{proof}

\begin{lemma}\label{lem-polynomial-width}
The parameters in \eqref{eq-global-parameters} satisfy
$r^2\le W\le352r^9$ and $r^2\le\rho<r^2+1$.
\end{lemma}
\begin{proof}
Cancellation of factorials gives
\[
 \beta=\frac{r(r-1)(r-2)^2}{2(64r^8+1)}.
\]
For $r\ge6$, the inequalities $r-1\ge5r/6$ and
$r-2\ge2r/3$ imply $\beta\ge1/(351r^4)$.
Also $0<\beta<1$. Thus
\[
 r^2\le W=r^2A<\frac{r^5}{\beta}+r^2
 \le351r^9+r^2\le352r^9.
\]
The defining ceiling for $A$ gives
$r^3\le A\beta<r^3+\beta$. Division by $r$ yields the bounds on
$\rho$.
\end{proof}

\section{From source contours to covering words}
\label{sec-assembly}

The numerical recurrence used later concerns finite lists of rows and
circles. We describe those lists and prove that an extension may use any
positive number of new letters. The one-letter transport lemmas in the
supplied assembly library apply repeatedly to give this construction.

\subsection{Rows, closed contours, and insertion}

Write $\operatorname{rot}(b,j)$ for the word obtained by rotating a nonempty
word $b$ to the left through $j$ positions. A \emph{cyclic class} consists
of all rotations of a word with distinct letters. If $x$ is absent from
$b$, write $\operatorname{ins}(b,x,j)$ for insertion of $x$ immediately
before position $j$, where positions start at zero and position $|b|$
means insertion at the end. Thus
$\operatorname{ins}(abc,x,1)=axbc$ and
$\operatorname{ins}(abc,x,3)=abcx$.

A \emph{row} is a triple $(b,s,v)$, where $b$ has distinct letters, the
letter $s$ does not occur in $b$, and $1\le v\le |b|$. The word $b$ is its
\emph{base}, the letter $s$ is its \emph{satellite}, and its \emph{charge}
is $|b|-v$. Its head and tail are the following literal words:
\[
 \operatorname{head}(b,s,v)=b[0:|b|-2],\qquad
 \operatorname{tail}(b,s,v)
   =\operatorname{rot}(b,v+1)[0:|b|-2].
\]
Here $b[i:j]$ means the consecutive letters in positions $i$ through
$j-1$. A row is \emph{full} if its charge is zero, and \emph{short} if
its charge is two. Except during the final completion, all rows below are
full or short. All rows in a given contour have the same satellite and the
same base alphabet.

Two rows are \emph{compatible} if the first tail equals the second head
as a word. A \emph{closed contour} is a nonempty finite list of rows in
which each row is compatible with the next, including the last and first
rows. Repeated occurrences remain separate entries of the list. For
example, the full rows with bases $abcd$, $bcad$, and $cabd$, all with one
fixed satellite outside $\{a,b,c,d\}$, form a closed contour. Their heads
are $ab,bc,ca$, respectively, and their tails are $bc,ca,ab$.

For a base of length $h+1$, an \emph{insertion port} is an integer in
$\{0,\ldots,h-1\}$. These are the $h$ insertion positions in a head of
length $h-1$. Across a full row the port changes from $p$ to $p-1$ modulo
$h$, and across a short row it changes from $p$ to $p+1$ modulo $h$.
Consequently the port displacement of a row is its charge minus one.
If a contour has $N$ rows and total charge $q$, its displacement is
$\Delta=q-N$.

We recall the literal transport operation from the supplied assembly
library~\cite{pantone-core}. Let $z$ be a fresh ordinary letter. A full
row $(b,s,h+1)$ at port zero is replaced by the two full rows with bases
$\operatorname{ins}(b,z,0)$ and $\operatorname{ins}(b,z,h)$.
At a port $1\le p<h$ it is replaced by the single full row with base
$\operatorname{ins}(b,z,p)$. A short row at a port $p<h-1$ is replaced
by the single short row with base $\operatorname{ins}(b,z,p)$.
At port $h-1$ it is replaced by two short rows with bases
$\operatorname{ins}(b,z,h-1)$ and
$\operatorname{rot}(b,h)z$. The satellite is unchanged in all four cases.
A transported contour is formed by concatenating these local paths and
propagating the port between successive source rows.

Each local path starts at the old head with $z$ inserted at the input
port and ends at the old tail with $z$ inserted at the output port.
This follows directly from the displayed definitions of head and tail.
For example, when $b=aUb$ is a full base and the input port is zero,
the two new bases are $zaUb$ and $aUzb$. Their common interface is the
literal word $aU$. These endpoint equations explain why compatible
source rows give compatible transported paths.

Across all $h$ input ports, each source row gives exactly $h+1$ new row
occurrences. All those occurrences have its charge. Their base cyclic
classes are precisely the classes obtained by inserting $z$ in all cyclic
gaps of the source base. The endpoint, inventory, and coverage statements
are the local transport propositions imported from the supplied
library~\cite{pantone-core}. The formulas above specify the operation to
which those propositions apply.

\begin{lemma}\label{lem-assembly-orbits}
Let a closed contour have bases of length $h+1$, where $h\ge2$, and put
$g=\gcd(h,|\Delta|)$ and $\ell=h/g$. For each $0\le e<g$, transport
$\ell$ consecutive copies of the contour from initial port $e$.
The resulting $g$ contours are closed. Their displacements are divisible
by $h+1$, and their combined row and charge counts are $(h+1)N$ and
$(h+1)q$, respectively. Every source-row and source-port pair occurs once
in their transport.
\end{lemma}

\begin{proof}
The ports at the beginnings of the successive source laps in component
$e$ are $e+t\Delta$ modulo $h$, for $0\le t<\ell$. Consider the map
\[
 (e,t)\longmapsto e+t\Delta\pmod h,
 \qquad 0\le e<g,\quad 0\le t<h/g.
\]
Equality of two outputs first gives equality of $e$ modulo $g$, and thus
equality of $e$ in its specified range. Division of the remaining
congruence by $g$ then gives
$(t-t')\Delta/g\equiv0\pmod{h/g}$. The integers $\Delta/g$ and $h/g$
are coprime, so $t=t'$ in the specified range. The domain and codomain
both have $h$ elements, so the map is a bijection. If $\Delta=0$, this
argument reduces to $g=h$, $\ell=1$, and the identity enumeration of
the ports.

After $\ell$ laps the source port returns to $e$, since $h$ divides
$\ell\Delta$. The literal endpoint equations therefore close the
transported path. The winding proposition for transport gives displacement
$(h+1)\Delta/g$ for that path, which is divisible by $h+1$. This is the
repeated-contour case of the imported winding theorem; its divisibility
hypothesis is exactly $h\mid\ell\Delta$.

At any fixed source-row occurrence, the prefix port change is a translation
of the set of lap-start ports. The bijection above therefore shows that
all its input ports occur once across the $g$ components. The local
inventory now gives $h+1$ output rows, each with the source charge, for
every source-row occurrence. Summing proves both count formulas.
\end{proof}

\subsection{Explicit circles meeting every completion port}

We reserve a second fresh letter $f$ for final completion. For precision,
we describe the completion operation used in the imported joining theorem.
For a row $(b,s,v)$, its lift at position $j$ has base
$\operatorname{ins}(b,f,j)$, satellite $s$, and visible count
$v+\boldsymbol{1}_{j<v}$. Its repair row at position $j$ has base
$\operatorname{rot}(b,j)s$, satellite $f$, and visible count $|b|+1$.
For a full row the completion path is its full transport path with new
letter $f$. For a short row at port $p<h-1$ the path is its one lift at
$p$. At port $h-1$ the path is, in order, the lift at $h-1$, the repair
rows at $h-1$ and $h$, and the lift at $h$. As before, concatenate paths
and propagate the source port. This defines the \emph{completed contour}
from each tagged initial port.

A \emph{circle} is a word of distinct letters, considered together with
all its rotations. Its displayed representative and its position in a
list of circles are retained for counting. For a contour with bases of
length $h+1$, the auxiliary circles have length $h$. We say that a circle
\emph{safely meets} a tagged completion port if either it contains $f$
and one of its rotations is the head of a row in the completed contour,
or the source contour is the full chart of a rotation of that circle.
The \emph{full chart} of a word $d$ and a letter $a$ outside $d$ consists
of the full rows with bases $\operatorname{rot}(d,j)a$, one for each
$0\le j<|d|$. This definition includes the circle $d$ used when a pending
class is closed. A list of circles is \emph{safe} for a contour if every
tagged initial completion port is safely met by at least one circle in
the list.

One useful sufficient witness for the first kind of meeting consists of
an actual source-row occurrence, the completion port propagated to that
occurrence, and a rotation identity
\[
 \operatorname{rot}(C,t)
 =\operatorname{ins}(\operatorname{head}(b,s,v),f,p).
\]
The first row of the local completion path has this head. Thus the
identity witnesses an actual vertex of the completed contour. Both the
occurrence and the propagated port matter; an equality of unlabeled
cyclic gaps is insufficient.

\begin{lemma}\label{lem-assembly-prefix}
Suppose a source path of port modulus $h$ starts at port $e$, where
$0\le e<h$. If its displacement is $S$, the displacement of its
transported path is
\[
 S+\left\lfloor\frac{e+S}{h}\right\rfloor.
\]
Immediately before the second transported row at a full source row whose
input port is zero, the transported prefix displacement is congruent to
$-e-1$ modulo $h+1$.
\end{lemma}

\begin{proof}
The formula is zero for the empty path. A full step subtracts one from
$S$. The floor decreases by one precisely when the old residue of $e+S$
is zero. Thus the claimed expression decreases by two at port zero and
by one at every other port. These are exactly the numbers of full rows
in the two possible full transport paths. A short step adds one to $S$.
The floor increases by one precisely at port $h-1$, so the expression
increases by two there and by one elsewhere. These are exactly the
displacements of the short transport paths. Induction proves the formula.
The floor is the Euclidean floor, also when $e+S$ is negative.

At a full port-zero occurrence, write $e+S=ht$. The first of its two
transported rows is full. Just before the second row, the displacement
is therefore $S+t-1=(h+1)t-e-1$, which gives the congruence.
\end{proof}

\begin{lemma}\label{lem-assembly-connectors}
In the setting of Lemma~\ref{lem-assembly-orbits}, suppose the source
contour has a consecutive run of $R$ full rows with bases
\[
 \operatorname{rot}(H,j)a,\qquad 0\le j<R,
\]
where $|H|=h$, $a$ is absent from $H$, and $R\le h$.
If $2g\lfloor R/g\rfloor\ge h$, the transported contours admit a safe
list of at most $2g^2$ circles. Each circle is supported on the source
alphabet together with $f$ and omits the satellite.
\end{lemma}

\begin{proof}
Put $L=\lfloor R/g\rfloor$, and let $S_0$ be the displacement before
the run begins. Fix a component with initial source port $e<g$, and
let $a_e$ be the residue of $e+S_0$ in $\{0,\ldots,g-1\}$.
The full source rows at phases
\[
 j=a_e+gu,\qquad 0\le u<L,
\]
all lie in the run. Indeed, $a_e+g(L-1)\le gL-1<R$.
Each is reached at port zero in exactly one of the $\ell$ laps. To
see this, its input port in lap $t$ is
$e+t\Delta+S_0-j$ modulo $h$. The congruence modulo $g$ holds by the
choice of $a_e$, and after division by $g$ the remaining congruence has
one solution for $t$ modulo $\ell$.

Write $H_e=\operatorname{rot}(H,a_e)$. For each $0\le t<g$, use the
two circles
\[
 C^-_t=\operatorname{ins}(H_e,f,g(L-1)+t),\qquad
 C^+_t=\operatorname{ins}(H_e,f,h-g+t).
\]
Their insertion positions lie between zero and $h-1$, because
$gL\le R\le h$. Both circles have length $h+1$, contain $f$, and
have distinct letters. At the second transported row of phase $a_e+gu$,
the head is $\operatorname{rot}(H_e,gu)$. In particular the transport
letter $z$ is absent from this head.

For $0\le j\le b<h$, splitting a word at positions $j$ and $b$ gives
\begin{align*}
 \operatorname{rot}(\operatorname{ins}(H_e,f,b),j)
   &=\operatorname{ins}(\operatorname{rot}(H_e,j),f,b-j),\\
 \operatorname{rot}(\operatorname{ins}(H_e,f,b),b+1)
   &=\operatorname{ins}(\operatorname{rot}(H_e,b),f,h).
\end{align*}
The second identity is a separate endpoint identity. It is needed because
the literal ports zero and $h$ of the advanced contour are distinct.

Given a port $gv+t$ with $0\le v<L$ and $0\le t<g$, choose
$u=L-1-v$ in the first identity for $C^-_t$. Thus the low circles meet
local ports $0,\ldots,gL-1$. The high circles, with $u$ varying through
$0,\ldots,L-1$, meet local ports $h-gL,\ldots,h-1$. For $C^-_0$ at
$u=L-1$, the second identity additionally meets port $h$.
The hypothesis $2gL\ge h$ says that the first two intervals cover
$0,\ldots,h-1$.

Finally fix a tagged initial completion port $J$ of this advanced
component. Lemma~\ref{lem-assembly-prefix} shows that its local port at
every selected second row is $J-e-1$ modulo $h+1$. Choose the local port
representative in $\{0,\ldots,h\}$ and use the corresponding circle
and occurrence above. The resulting literal identity proves safe
incidence from $J$. There are $2g$ circles for each of $g$ components,
so the total is $2g^2$. Coincident circles may remain as separate list
occurrences. Every circle uses only $H_e$ and $f$, which proves the
support statement.
\end{proof}

The literal hypothesis on the full run is necessary. Compatibility alone
allows a different full successor. For example, the three full bases
$012,120,201$ form a compatible contour, but their first two-letter
prefixes are not rotations of one fixed two-letter word. The positive-slot
construction must supply a run from an untouched normal batch, where
successive bases have precisely the form required in
Lemma~\ref{lem-assembly-connectors}.

\subsection{The finite universe above an old cyclic class}

Fix a word $d$ on $m\ge2$ distinct old letters and distinct new letters
$v_0,\ldots,v_{W-1}$. A cyclic class on these letters is \emph{above}
$d$ if deletion of all the $v_i$ leaves the cyclic class of $d$.
Define lists $U_i(d)$ recursively. Start with $U_0(d)=[d]$. For each
$b$ in $U_i(d)$, place
\[
 \operatorname{rot}(b,j)v_i,\qquad 0\le j<m+i,
\]
in $U_{i+1}(d)$, retaining the order of the parent list and the indicated
order of $j$. For example, from $d=ab$ and one new letter $x$ this gives
$U_1(d)=[abx,bax]$.

\begin{lemma}\label{lem-assembly-universe}
The list $U_W(d)$ contains exactly one representative of each cyclic
class above $d$. Its length is
\[
 D_{m,W}=\prod_{i=0}^{W-1}(m+i).
\]
If $B$ source-row occurrences have pairwise distinct base cyclic classes
above $d$, their complement in $U_W(d)$ is an actual list of length
$D_{m,W}-B$.
\end{lemma}

\begin{proof}
We prove the first assertion by induction on $i$. It holds for $i=0$.
Given a class above $d$ on the first $i+1$ new letters, delete $v_i$.
The remaining class has a unique representative $b$ in $U_i(d)$ by
induction. Rotate the larger word so that its unique occurrence of $v_i$
is last. The preceding word is a rotation of $b$, which gives existence
of a representative in the displayed list.

For uniqueness, equality of two larger cyclic classes implies equality
of their deleted parent classes. Thus their parent representatives agree.
Because $v_i$ occurs once and is last in both displayed representatives,
cyclic equivalence forces equality of those representatives as literal
words. A word of $m+i$ distinct letters has $m+i$ different rotations,
so the two rotation indices agree as well. Each parent therefore
contributes exactly $m+i$ new representatives. Multiplication over the
steps gives the formula for $D_{m,W}$.

Each selected source occurrence has a unique representative in $U_W(d)$.
The distinctness assumption makes these representatives distinct. Delete
these $B$ entries from the finite list. The resulting complement has
length $D_{m,W}-B$, as asserted. This counting operation does not change
the literal bases in the source contours.
\end{proof}

\subsection{Extension by an arbitrary number of letters}

We now specify the intermediate objects. Reserve distinct letters $s$ and
$f$ for the satellite and final completion. A \emph{stage of old size}
$m$ consists of an old alphabet $A_0$ of size $m$, a distinguished
ordinary letter $z$ outside $A_0$, a list of contours on $z\cup A_0$,
a list of auxiliary circles of length $m$, and a list of pending cyclic
classes on $A_0$. The alphabets omit $s$ and $f$. Each contour has
displacement divisible by $m$. Every circle has distinct letters and is
supported on $\{f,z\}\cup A_0$, and the circle list is safe for every
contour. Finally every cyclic class on $z\cup A_0$ either occurs as a
row base or is obtained by inserting $z$ into a pending cyclic class.
Write $R,Q,c,P$ for the row count, total charge, circle count, and pending
count. Counts refer to list occurrences.

We use the following existing transport facts. Transporting every
component from every port through a fresh ordinary letter preserves
closure, displacement divisibility and coverage. It multiplies row and
charge counts by $m+1$. Extending each circle through all its $m$ cyclic
insertion positions gives $mc$ circles, preserves distinct letters and
support, and safely meets every transported component. These statements
are the general one-letter level and circle transport results in the
supplied assembly library~\cite{halfar-assembly}. They apply to arbitrary
stages satisfying the preceding conditions.

\begin{lemma}\label{lem-assembly-width}
Let a stage have old size $m\ge2$, counts $(R,Q,c,P)$, old alphabet
$A_0$, and distinguished letter $v_0$. Let $W\ge1$, and choose distinct
fresh ordinary letters $v_1,\ldots,v_{W-1},z'$. Suppose that, above
every pending old class $d$, there is a family of closed source contours
on $A_0\cup\{v_0,\ldots,v_{W-1}\}$ with $B$ pairwise distinct source
classes above $d$ and total charge $q$. Suppose its transport through
$z'$ has displacement divisible by $m+W$ and admits a safe supported
list of $c_{\mathrm{adv}}$ circles of length $m+W$.
Then there is a stage of old size $m+W$ and distinguished letter $z'$
with counts
\begin{align*}
 R'&=A_{m,W}R+(m+W)PB,\\
 Q'&=A_{m,W}Q+(m+W)Pq,\\
 c'&=D_{m,W}c+Pc_{\mathrm{adv}},\\
 P'&=P(D_{m,W}-B),
\end{align*}
where $A_{m,W}=\prod_{i=1}^{W}(m+i)$ and $D_{m,W}$ is as in
Lemma~\ref{lem-assembly-universe}. The transport of the source family
is required to include every source-row and source-port pair once.
\end{lemma}

\begin{proof}
Set $A'_0=A_0\cup\{v_0,\ldots,v_{W-1}\}$. First transport the
existing contours successively through
$v_1,\ldots,v_{W-1},z'$, and extend their circles through the same
sequence. There are exactly $W$ transports: the old distinguished letter
$v_0$ already belongs to every existing component. At each step the
letter being added is absent from the current ordinary alphabet and from
both reserved roles. The one-letter results therefore apply inductively.
They give closed contours on $z'\cup A'_0$, displacement divisible by
$m+W$, valid supported circles, and safety from every completion port.
They multiply the existing row and charge counts by
$(m+1)\cdots(m+W)=A_{m,W}$ and the circle count by
$m\cdots(m+W-1)=D_{m,W}$.

For every old pending occurrence $d$, append its transported source
contours and its supplied circles. Their structural hypotheses put them
on the same alphabet with the same port modulus as the transported old
components. Taking the union preserves closure and displacement
divisibility component by component. It also preserves circle validity
and support. A circle which met a component before taking the union is
still in the union, so safety is preserved. Replace $d$ in the pending
list by the actual complement list from
Lemma~\ref{lem-assembly-universe}. Every new pending representative is
a permutation of $A'_0$.

It remains to prove the coverage condition. Let $y$ be an arbitrary
cyclic class on $z'\cup A'_0$. Delete the new letters in the reverse
order $z',v_{W-1},\ldots,v_1$. This produces a cyclic class $x$ on
$v_0\cup A_0$. Apply the old coverage condition to $x$.

If $x$ is the class of an existing row, reinsert
$v_1,\ldots,v_{W-1},z'$ in the gaps determined by $y$. The one-letter
coverage proposition at each transport gives a final transported row
in class $y$.

Otherwise $x$ is above an old pending class $d$. Delete $z'$ from $y$
and call the resulting source class $t$. Deleting all $W$ source
controls from $t$ gives $d$, so $t$ has one representative in $U_W(d)$.
If that representative was selected, the source family has a row with
base class $t$. The exhaustive transport inventory and the local
insertion coverage supply a transported row in the class $y$. If it was
not selected, its representative lies in the new pending complement,
and $y$ is obtained from that representative by inserting $z'$. These
two alternatives establish the new coverage condition.

Each pending old occurrence contributes $(m+W)B$ advanced rows and
charge $(m+W)q$, by the source-row and source-port inventory. It
contributes exactly $c_{\mathrm{adv}}$ circles and $D_{m,W}-B$ pending
classes. Adding these contributions to the transported old objects gives
all four displayed counts.
\end{proof}

For the source family of Proposition~\ref{prop-global-counts}, each
surviving candidate contributes the closed contour obtained by its
specified modifications. Order the finite set of surviving candidates
and list these contours in that order. The local distinctness theorem
and the exclusion of collisions between candidates show that all their
source-row cyclic classes are distinct. Their row and charge counts are
the sums of the counts of those same contours. In particular, selecting
this list requires no additional choice of rows or cyclic classes.

Each contour has an untouched normal batch and thus a full run of the
form required in Lemma~\ref{lem-assembly-connectors}. The verification
in Proposition~\ref{prop-global-counts} gives the inequality for its
length. Lemmas~\ref{lem-assembly-orbits}
and~\ref{lem-assembly-connectors} consequently provide the transported
contours, their complete source-port inventory, and their safe circle
lists. Thus the family satisfies every source hypothesis of
Lemma~\ref{lem-assembly-width}.

A family constructed on one fixed old order can be used above another
old order by a single bijection of the full finite alphabet. Apply this
bijection to the old letters, all source controls, the satellite, the
transport letter, and the completion letter together. Rotation,
insertion, prefixes, and equality of words commute with this relabeling.
It therefore preserves literal compatibility and every incidence
witness, as well as the counts. Relabeling only part of a frame would
not justify those equalities.

\subsection{An imported family for the final extensions}

For the last part of the schedule, we use an existing family with three
source controls. We state both its geometric properties and the numerical
bounds that will be used. Put
\[
 \mathcal G_a=\{g\in\mathbb Z:0\le g<a,\quad
   g\equiv1\pmod3,\quad 2g+5\le a\},
 \qquad
 \nu_a=
 \begin{cases}2,&a\text{ odd},\\8,&a\text{ even}.\end{cases}
\]

\begin{proposition}\label{prop-assembly-dense}
For every old cyclic order on $a\ge7$ letters, three distinct source
controls, a fresh transport letter, and two further distinct reserved
letters, there is a source family satisfying the hypotheses of
Lemma~\ref{lem-assembly-width} with width three. Its counts are
\[
 B_a=2|\mathcal G_a|a(a+12),\qquad
 q_a=16|\mathcal G_a|a,\qquad
 c_{\mathrm{adv},a}=(a+2)\nu_a.
\]
Writing $D_{a,3}=a(a+1)(a+2)$, these counts satisfy
$B_a\le D_{a,3}$ and $\nu_a\le8$. For $a\ge49$, they also satisfy
\[
 D_{a,3}\le4B_a,\qquad
 a(q_a+c_{\mathrm{adv},a})\le8B_a.
\]
\end{proposition}

This is the family \texttt{UniR.rFam} in the supplied assembly
library~\cite{halfar-assembly}. Its construction is valid for an
arbitrary full frame of distinct letters satisfying the two reserved-letter
exclusions. The declarations \texttt{rData\_ok} and
\texttt{rData\_safe} give literal source validity, distinct source classes,
and safe incidence from every required port. The constructor
\texttt{Uniform.famS} supplies the transported components, their
displacement divisibility, and their finite complement. Thus this
proposition gives actual families and circle lists for
Lemma~\ref{lem-assembly-width}.

The count formulas and the bounds $B_a\le D_{a,3}$, $\nu_a\le8$, and
$D_{a,3}\le4B_a$ are respectively \texttt{rowsR\_eq}, \texttt{chargeR\_eq},
\texttt{rowsR\_le}, \texttt{ncycR\_le}, and
\texttt{rows\_quarter} in \texttt{UniR/RRate.lean}.
The cost inequality is \texttt{cost\_le8} in
\texttt{UniR/RRate8.lean}. The count of advanced circles is
$(a+2)\nu_a$, since each of the original $\nu_a$ circles has $a+2$
extensions at transport. Whole-frame relabeling as described above
transfers these statements to the chosen old cyclic order and ordinary
letters. No parity restriction remains, and the family can extend an
already existing stage. It need not start a new construction.

\subsection{Initialization, padding, and the final word}

At old size $m$, choose no contours and no circles, and take every cyclic
class on the old alphabet to be pending. There are $(m-1)!$ such classes:
fix one old letter first and permute the remaining $m-1$ letters. Every
class on the old alphabet together with the distinguished letter lies
above one of them, by deletion of that letter. Thus this is a stage
with counts $(0,0,0,(m-1)!)$.

The identity
\[
 (m+W)D_{m,W}=mA_{m,W}
\]
and Lemma~\ref{lem-assembly-width} show that
\[
 R'+(m+W)P'=A_{m,W}(R+mP).
\]
Consequently every stage constructed from this initialization satisfies
$R+mP=m!$. Taking $W=1$ and an empty selected family gives a valid
one-letter extension with
\[
 (R',Q',c',P')=((m+1)R,(m+1)Q,mc,mP).
\]
Such extensions permit adjustment of the final alphabet size without a
congruence restriction.

At a final old size $h$, replace each pending representative $d$ by its
full chart with the current distinguished letter, and add the circle
$d$. The chart is closed, has $h$ rows, has zero charge, and has
displacement $-h$. It covers every class above $d$ and is safely met by
its own circle according to the definition. Thus the resulting closed
level has $h!$ rows, charge $Q$, and $c+P$ circles, and covers every base
cyclic class.

We use the following precise consequence of the imported word assembly
theorem~\cite{halfar-assembly}. Suppose $h\ge8$. A closed level on an ordinary alphabet
of size $h+1$, with $h!$ rows, charge $Q$, and a valid supported safe
list of $C$ circles, gives a covering word on $n=h+3$ letters of length
at most
\[
 F_3(n)+Q+hC+(a-2)C
 +(h+1-a)\min\{C,(n-1)^{\underline{h+1-a}}\},
 \qquad 2\le a\le h+1.
\]
The notation $x^{\underline{k}}$ means $x(x-1)\cdots(x-k+1)$.
The completed alphabet, consisting of the ordinary letters, satellite,
and completion letter, is relabeled bijectively to $\{0,\ldots,n-1\}$,
with the satellite labeled $9$ and the completion letter labeled $10$.
These labels are available because $n=h+3\ge11$. All other labels can
be assigned bijectively to the ordinary letters. The relabeling preserves
contiguous occurrences and meets the alphabet hypothesis of the imported
declaration \texttt{Closed.word}. The theorem supplies
an actual word containing every permutation, rather than an estimate
for an abstract count.

Taking $a=2$, $C=c+P$, and bounding the minimum by $C$ gives
\begin{equation}
 |w|\le F_3(n)+Q+(2h-1)(c+P).
 \label{eq-assembly-word}
\end{equation}
This is the completion estimate used in the finite transfer argument.
The source families of Proposition~\ref{prop-global-counts}, followed by
the families of Proposition~\ref{prop-assembly-dense}, maintain all of
these hypotheses through the extensions above.

% The abstract transfer has explicit construction hypotheses. Its final
% corollaries instantiate them with the families in the preceding sections.
\section{A finite transfer estimate}
\label{sec-finite-transfer}

We now estimate the length supplied by the assembly construction. The
starting size will be close to $n$, and most extensions will use the
families of Proposition~\ref{prop-global-counts}. A final sequence of
extensions of size three reduces the pending fraction far enough for
completion. We first give the transfer with explicit hypotheses and then
verify them using the constructions in the preceding sections.

\subsection{The construction statements needed for the estimate}

For a stage of old size $m\ge2$, retain the counts $R,Q,c,P$ from
Section~\ref{sec-assembly}. Define the rational numbers
\[
 y=\frac{Q}{m!},\qquad
 z=\frac{c}{(m-1)!},\qquad
 p=\frac{P}{(m-1)!},\qquad X=y+z.
\]
There are $(m-1)!$ cyclic classes on $m$ letters, so $p$ is the fraction
of classes still awaiting an assignment. Normalizing the other counts by
the indicated factorials will make inherited contributions unchanged
under an extension.

An extension from size $m$ to size $m+v$, where $v$ is a positive integer,
uses nonnegative integer counts $B,q,c_{\mathrm{adv}}$. Put
\[
 D_{m,v}=\prod_{i=0}^{v-1}(m+i),\qquad
 A_{m,v}=\prod_{i=1}^{v}(m+i).
\]
Lemma~\ref{lem-assembly-width} gives the bounds and update rules
\begin{align}
 0&\le B\le D_{m,v}, \label{eq-transfer-source-bound}\\
 R'&=A_{m,v}R+(m+v)PB,\notag\\
 Q'&=A_{m,v}Q+(m+v)Pq,\notag\\
 c'&=D_{m,v}c+Pc_{\mathrm{adv}},\notag\\
 P'&=P(D_{m,v}-B). \label{eq-transfer-integer-update}
\end{align}
The distinctness hypothesis in Lemma~\ref{lem-assembly-universe} is
required for \eqref{eq-transfer-source-bound} and for the complement
count in the last line.

Assume that the construction can start at any size $m_0\ge11$ with
\[
 (R,Q,c,P)=(0,0,0,(m_0-1)!).
\]
Use also the empty extensions from Section~\ref{sec-assembly}, which have
size one and $B=q=c_{\mathrm{adv}}=0$. At the final intermediate size
$h=n-3$, bounding the minimum in \eqref{eq-assembly-word} by $c+P$
gives an actual superpermutation $w$ with
\begin{equation}
 |w|\le F_3(n)+Q+(2h-1)(c+P).
 \label{eq-transfer-word-input}
\end{equation}
In the conditional results below, the initialization, the required
extensions, and completion to an actual superpermutation are all hypotheses. Equations
for counts without those construction statements are insufficient.

\begin{lemma}\label{lem-transfer-normalization}
The rules \eqref{eq-transfer-integer-update} preserve the identity
$R+mP=m!$. With
\[
 f=\frac{B}{D_{m,v}},\qquad
 g=\frac{q+c_{\mathrm{adv}}}{D_{m,v}},
\]
their normalized form is
\begin{equation}
 p'=p(1-f),\qquad
 y'=y+p\frac{q}{D_{m,v}},\qquad
 z'=z+p\frac{c_{\mathrm{adv}}}{D_{m,v}},\qquad
 X'=X+pg.
 \label{eq-transfer-normalized-update}
\end{equation}
In particular, $0\le p\le1$ throughout the construction, and empty
extensions preserve $y,z,p$.
\end{lemma}

\begin{proof}
Cancellation in the products gives
\[
 (m+v)D_{m,v}=mA_{m,v},\qquad
 (m+v)!=A_{m,v}m!,\qquad
 (m+v-1)!=D_{m,v}(m-1)!.
\]
Consequently,
\[
 R'+(m+v)P'
 =A_{m,v}R+(m+v)PD_{m,v}
 =A_{m,v}(R+mP).
\]
This is $(m+v)!$ when $R+mP=m!$. The initial counts satisfy that identity.
Dividing the updates for $Q',c',P'$ by the appropriate new factorials
proves \eqref{eq-transfer-normalized-update}. Initially $p=1$.
The inequality $0\le f\le1$ therefore proves $0\le p\le1$ by induction.
For an empty extension, $f=g=q=c_{\mathrm{adv}}=0$, so the normalized
quantities do not change.
\end{proof}

\subsection{Two finite inequalities}

The main estimate uses two short induction arguments. They apply to
rational numbers and do not require an approximation to a product.

\begin{lemma}\label{lem-transfer-pending}
Let $f_0,\ldots,f_{J-1}$ lie in $[0,1]$. Set $p_0=1$ and
$p_{i+1}=p_i(1-f_i)$. Then
\[
 p_J\left(1+\sum_{i=0}^{J-1}f_i\right)\le1.
\]
\end{lemma}

\begin{proof}
Write $s_j=\sum_{i<j}f_i$. All the numbers $p_j,s_j$ are nonnegative.
Using $s_{j+1}=s_j+f_j$, we obtain
\begin{align*}
 p_{j+1}(1+s_{j+1})
 &=p_j(1-f_j)(1+s_j+f_j)\\
 &=p_j(1+s_j)-p_jf_j(s_j+f_j)\\
 &\le p_j(1+s_j).
\end{align*}
The initial value is $p_0(1+s_0)=1$. Induction proves the claim.
\end{proof}

For example, if $f_0=f_1=1/2$, then $p_2=1/4$, and the left side of
the lemma is $1/2$. The estimate remains valid if a selected density is
one, in which case all subsequent pending mass is zero.

\begin{lemma}\label{lem-transfer-cost}
Suppose $p_{i+1}=p_i(1-f_i)$ and $X_{i+1}=X_i+p_i g_i$, with
$p_i\ge0$. If $t g_i\le Kf_i$ for every $i<J$, where $t>0$, then
\[
 t(X_J-X_0)\le K(p_0-p_J).
\]
\end{lemma}

\begin{proof}
For each $i$,
\[
 t(X_{i+1}-X_i)=tp_i g_i
 \le Kp_i f_i=K(p_i-p_{i+1}).
\]
Summing cancels all intermediate $X_i$ and $p_i$ and gives the result.
\end{proof}

\subsection{The schedule and its cost}

Fix an integer $r\ge6$ divisible by six. Let $W$ be a positive multiple
of $r$, let $E>0$, let $\rho\ge r^2$, and put
\[
 c_r=\frac{2(r-1)}{r-2}.
\]
At every multiple $m$ of $r$ with $m\ge4E$, assume that the construction
provides an extension of size $W$ for which
\begin{equation}
 f_m=\frac{\rho W}{m}(1+\theta_m),\qquad
 g_m=\frac{c_r\rho W}{m^2}(1+\varepsilon_m),\qquad
 |\theta_m|,|\varepsilon_m|\le\frac{E}{m},
 \label{eq-transfer-sparse-input}
\end{equation}
and
\begin{equation}
 c_{\mathrm{adv},m}\le\frac{r-1}{m}q_m.
 \label{eq-transfer-circle-input}
\end{equation}
The final extensions come from Proposition~\ref{prop-assembly-dense}.
For an extension of size three at any $a\ge49$, its inequalities and
Lemma~\ref{lem-transfer-normalization} give
\begin{equation}
 p'\le\frac34p,\qquad
 X'-X\le\frac8a(p-p'),\qquad
 z'-z\le\frac{8p}{a(a+1)}.
 \label{eq-transfer-dense-input}
\end{equation}
The extensions of size $W$ select a small proportion of the available
classes. The extensions of size three remove at least one quarter of
the pending mass. We use the latter only at the end.

Here is the derivation of \eqref{eq-transfer-dense-input}. Its selected
fraction is $f=B_a/D_{a,3}\ge1/4$, so
$p'=p(1-f)\le3p/4$. The cost estimate of
Proposition~\ref{prop-assembly-dense} gives
\[
 X'-X=p\frac{q_a+c_{\mathrm{adv},a}}{D_{a,3}}
 \le\frac8a\frac{pB_a}{D_{a,3}}
 =\frac8a(p-p').
\]
Finally, its circle count gives
\[
 z'-z=p\frac{(a+2)\nu_a}{a(a+1)(a+2)}
 =\frac{p\nu_a}{a(a+1)}\le\frac{8p}{a(a+1)}.
\]
These are updates of the existing stage. In particular, the final
extensions retain the rows and safe circles obtained during the sparse
part of the construction.

For a positive integer $n$, let $\ell$ be the least nonnegative integer
such that $n<2^\ell$. Define
\[
 K=12\ell,\qquad \tau=3+36\ell.
\]
For example, $\ell=7$ when $n=64$, because the upper inequality in the
definition is strict.

\begin{theorem}\label{thm-transfer-finite}
Assume the construction and completion statements above. If
\begin{equation}
 n\ge8rE,\qquad n\ge96r^2,\qquad
 r+W+\tau\le\frac{n}{2r},
 \label{eq-transfer-room}
\end{equation}
then there is a superpermutation $w$ on $n$ letters such that
\[
 |w|\le F_3(n)+\left(2+\frac{74}{r}\right)(n-4)!.
\]
\end{theorem}

\begin{proof}
Put $\alpha=(r-1)/r$, and let $m_0$ be the least multiple of $r$
that is at least $\alpha n$. Thus
\[
 \alpha n\le m_0<\alpha n+r.
\]
Set $a=n-\tau$ and
\[
 J=\left\lfloor\frac{a-m_0}{W}\right\rfloor,
 \qquad b=m_0+JW.
\]
Since $\lfloor x\rfloor>x-1$, the last condition in
\eqref{eq-transfer-room} gives
\[
 J>\frac{n/r-r-\tau-W}{W}
 \ge\frac{n}{2rW}>0.
\]
In particular $J\ge1$. Starting at $m_0$, use $J$ extensions of size
$W$, then $a-b$ empty extensions, and finally $K$ extensions of size
three. The floor defining $J$ gives $0\le a-b<W$. The final size is
$a+3K=n-3$, as required for completion.

The initial size satisfies $m_0\ge n/2\ge48r^2\ge11$. Each start
$m_i=m_0+iW$ of an extension of size $W$ is a multiple of
$r$ and lies between $n/2$ and $n$. Since $n\ge8rE$, it is at least
$4E$. Also $a\ge n/2$, because $\tau\le n/(2r)$, and
$a\ge n/2\ge48r^2>49$. Thus every extension used in this schedule
is within its stated range. The last extension of size three starts at
$n-6$.

Write
\[
 \delta=\frac{2E}{n}\le\frac1{4r}.
\]
The two relative errors at each $m_i$ have absolute value at most
$\delta$. In particular $f_{m_i}>0$. From
\eqref{eq-transfer-sparse-input},
\[
 \frac{ng_{m_i}}{f_{m_i}}
 \le\frac{c_r}{\alpha}\frac{1+\delta}{1-\delta}.
\]
Here
\[
 \frac{c_r}{\alpha}=\frac{2r}{r-2}
 =2+\frac4{r-2}\le2+\frac6r,
\]
because $r\ge6$. Since $0\le\delta\le1/2$,
\[
 \frac{1+\delta}{1-\delta}
 =1+\frac{2\delta}{1-\delta}
 \le1+4\delta\le1+\frac1r.
\]
Finally,
\[
 \left(2+\frac6r\right)\left(1+\frac1r\right)
 =2+\frac8r+\frac6{r^2}
 \le2+\frac9r.
\]
Let $p_s,X_s,y_s,z_s$ be the normalized values after the extensions of
size $W$. Lemma~\ref{lem-transfer-cost}, with $t=n$, now gives
\begin{equation}
 nX_s\le\left(2+\frac9r\right)(1-p_s)
 \le2+\frac9r.
 \label{eq-transfer-sparse-cost}
\end{equation}

We next bound $p_s$. Each selected density satisfies
\[
 f_{m_i}\ge\frac{\rho W}{n}(1-\delta)
 \ge\frac{\rho W}{2n}.
\]
The lower bound on $J$ therefore yields
\[
 \sum_{i<J} f_{m_i}
 \ge\frac{n}{2rW}\frac{\rho W}{2n}
 =\frac{\rho}{4r}\ge\frac r4.
\]
Lemma~\ref{lem-transfer-pending} proves
\begin{equation}
 p_s\le\frac4{r+4}\le\frac4r.
 \label{eq-transfer-pending-small}
\end{equation}
The intervening empty extensions preserve these normalized quantities.

All the final extensions start at sizes at least $a$. Their cost bound
in \eqref{eq-transfer-dense-input} telescopes to
\[
 X_{\mathrm{end}}-X_s
 \le\frac8a(p_s-p_{\mathrm{end}})
 \le\frac{8p_s}{a}.
\]
Using $n/a\le2$, \eqref{eq-transfer-sparse-cost}, and
\eqref{eq-transfer-pending-small}, we obtain
\begin{equation}
 nX_{\mathrm{end}}\le2+\frac{73}{r}.
 \label{eq-transfer-total-cost}
\end{equation}

The completion estimate also involves $z_{\mathrm{end}}$. During the
extensions of size $W$, \eqref{eq-transfer-circle-input} gives
\[
 z_s\le\frac{r-1}{m_0}y_s.
\]
Their local cost satisfies $m_i g_{m_i}\le8f_{m_i}$, since
$m_i\le n$ and $2+9/r\le8$. Applying
Lemma~\ref{lem-transfer-cost} with $t=m_0$ gives
$y_s\le X_s\le8/m_0$. Hence
\[
 z_s\le\frac{8(r-1)}{m_0^2}.
\]
For the final extensions, write $p_j$ for the pending mass before the
$j$th one. The contraction $p_{j+1}\le3p_j/4$ implies
$p_j\le4(p_j-p_{j+1})$. Summing gives
$\sum_{j<K}p_j\le4(p_s-p_{\mathrm{end}})\le4$.
The last inequality in \eqref{eq-transfer-dense-input} consequently gives
\[
 z_{\mathrm{end}}-z_s
 \le\frac{32}{a(a+1)}.
\]
Both $m_0$ and $a$ are at least $n/2$, so
\begin{equation}
 z_{\mathrm{end}}\le\frac{32r+96}{n^2}.
 \label{eq-transfer-final-circles}
\end{equation}

There are $12\ell$ final extensions. Since $p_s\le1$ and
$(3/4)^4=81/256<1/2$, their contraction gives
\begin{equation}
 p_{\mathrm{end}}
 \le(3/4)^{12\ell}<2^{-3\ell}<n^{-3}.
 \label{eq-transfer-final-pending}
\end{equation}
The strict inequalities are valid because $\ell\ge1$ and $n<2^\ell$.

Put $h=n-3$. Divide \eqref{eq-transfer-word-input} by $(n-4)!=(h-1)!$.
The normalized excess is at most
\begin{align*}
 h y_{\mathrm{end}}+(2h-1)(z_{\mathrm{end}}+p_{\mathrm{end}})
 &=hX_{\mathrm{end}}+(h-1)z_{\mathrm{end}}
   +(2h-1)p_{\mathrm{end}}\\
 &\le nX_{\mathrm{end}}+nz_{\mathrm{end}}+2np_{\mathrm{end}}.
\end{align*}
Equations \eqref{eq-transfer-total-cost}--\eqref{eq-transfer-final-pending}
bound this by
\[
 2+\frac{73}{r}+\frac{32r+96}{n}+\frac2{n^2}.
\]
For $r\ge6$, we have $32r+96\le48r$. The assumption $n\ge96r^2$
therefore bounds the penultimate term by $1/(2r)$. It also implies
$n^2\ge4r$, so the last term is at most $1/(2r)$. Their sum is at most
$1/r$, which proves the theorem.
\end{proof}

\subsection{One explicit threshold}

Choose the parameters in \eqref{eq-global-parameters}. By
Lemma~\ref{lem-polynomial-width}, they satisfy
$r^2\le W\le352r^9$ and $r^2\le\rho<r^2+1$. We need one further
consequence of those bounds.

\begin{lemma}\label{lem-transfer-parameter-bounds}
The integer $W$ is divisible by $r$, and
$E\le10^{12}352^4r^{58}$.
\end{lemma}

\begin{proof}
The identity $W=r^2A$ proves divisibility. The upper bound for $W$ gives
\[
 E=10^{12}r^{22}W^4
 \le10^{12}r^{22}(352r^9)^4
 =10^{12}352^4r^{58}.
\]
\end{proof}

\begin{lemma}\label{lem-transfer-binary-room}
For these parameters, the condition $n\ge8rE$ implies every condition
in \eqref{eq-transfer-room}.
\end{lemma}

\begin{proof}
Since $W\ge r^2\ge1$, we have $E\ge W$ and
$E\ge10^{12}r^{22}\ge2592r$. Consequently
\[
 n\ge8rE\ge(144r)^2.
\]
This proves $n\ge64$ and $n\ge96r^2$.

Recall that $\ell$ is the least integer with $n<2^\ell$.
Thus $\ell\ge7$ and $2^{\ell-1}\le n$. For every integer $t\ge7$,
\[
 t^2\le2^{t-1}.
\]
Indeed, the initial case is $49\le64$, and
$(t+1)^2\le2t^2$ for $t\ge7$, which proves the induction step.
It follows that $\ell^2\le n$. If $\ell\le144r$, then
$144r\ell\le(144r)^2\le n$. If $\ell\ge144r$, then
$144r\ell\le\ell^2\le n$. In either case,
\[
 36\ell\le\frac{n}{4r}.
\]
Finally, $W\ge r^2\ge r+3$, and therefore
\[
 r+W+3\le2W\le2E\le\frac{n}{4r}.
\]
Adding the last two inequalities gives
$r+W+\tau\le n/(2r)$, as required.
\end{proof}

\begin{corollary}\label{cor-transfer-polynomial-threshold}
For every integer $r\ge6$ divisible by six and every integer
$n\ge Cr^{61}$, there is an actual superpermutation $w$ satisfying
\[
 r|w|\le rF_3(n)+(2r+74)(n-4)!.
\]
\end{corollary}

\begin{proof}
Use the parameters of \eqref{eq-global-parameters}. At each multiple
$m$ of $r$ with $m\ge4E$, Proposition~\ref{prop-global-counts}
supplies actual source families above any chosen old cyclic order. Its
distinct source classes lie in the universe of
Lemma~\ref{lem-assembly-universe}, so $0\le B\le D_{m,W}$ and the
required complement exists. The same proposition supplies the literal
full run needed by Lemma~\ref{lem-assembly-connectors}. Together with
Lemma~\ref{lem-assembly-orbits}, this gives the safe supported circle
lists and transported components required by
Lemma~\ref{lem-assembly-width}.

Initialization and empty extensions are provided in
Section~\ref{sec-assembly}. The final extensions and their normalized
bounds are provided by Proposition~\ref{prop-assembly-dense} and
\eqref{eq-transfer-dense-input}. The schedule has final old size $n-3$,
and $n\ge Cr^{61}\ge11$, so \eqref{eq-assembly-word} gives its actual
superpermutation. This verifies all the construction and completion
hypotheses of Theorem~\ref{thm-transfer-finite}.

Lemma~\ref{lem-transfer-parameter-bounds} gives
\[
 8rE\le Cr^{59}\le Cr^{61}\le n.
\]
Lemma~\ref{lem-transfer-binary-room} supplies the other size conditions.
Apply Theorem~\ref{thm-transfer-finite} and multiply its conclusion by
the positive integer $r$.
\end{proof}

\subsection{The exponent in the asymptotic statement}

\begin{corollary}\label{cor-transfer-asymptotic}
\[
 L(n)=F_3(n)+2(n-4)!+
 O\!\left(\frac{(n-4)!}{n^{1/61}}\right).
\]
\end{corollary}

\begin{proof}
For $n\ge C6^{61}$, choose the largest positive multiple $r$ of six
with $Cr^{61}\le n$. Such an $r$ exists, and maximality gives
\[
 Cr^{61}\le n<C(r+6)^{61}\le C(2r)^{61}.
\]
The last inequality uses $r\ge6$. Taking positive sixty-first roots
now gives
\[
 \frac1r<2C^{1/61}n^{-1/61}.
\]
The word from Corollary~\ref{cor-transfer-polynomial-threshold} yields
\[
 \frac{L(n)-F_3(n)}{(n-4)!}-2
 \le\frac{74}{r}
 <148C^{1/61}n^{-1/61}.
\]
On the other side, use $(n-4)!=(n-4)(n-5)!$ to divide the imported lower bound
\eqref{eq-main-lower} by $(n-4)!$. Since $n\ge8$, this gives
\[
 2-\frac{L(n)-F_3(n)}{(n-4)!}
 \le\frac2{n-4}\le\frac4n\le4n^{-1/61}.
\]
Both sides are bounded by a constant times $n^{-1/61}$. Multiplying
by $(n-4)!$ proves the stated remainder.
\end{proof}


\section{Formal verification and scope}

The accompanying formalization proves Theorem~\ref{thm-main-finite} and
Corollary~\ref{cor-main-asymptotic}. It constructs the literal source rows,
their repairs and safe connector circles, and the globally distinct source
families. The exact counts of these objects supply the estimates used by
the arbitrary-width extension. The rational transfer and the original
word-assembly theorem then produce an actual covering word. Combining its
length bound with the rebuilt lower-bound proof gives the two-sided result.

The final declarations are \texttt{finite\_covering\_word} and
\texttt{shortest\_length\_asymptotic} in the namespace
\texttt{CoefficientTwoPrecision}.
They use the original coverage predicate and minimum-length definition.
Neither theorem assumes a source family, a numerical recurrence, or the
existence of a word with the stated bound. The accompanying proof map
matches the mathematical steps to their Lean declarations.

The formalization uses Lean 4.35.0. Its pinned Mathlib revision is
\begin{center}
\texttt{56a005855ea17e0b0958f73ec801ae140b0a1679}.
\end{center}
Build receipts record the checked source hashes, dependency versions,
compiler results, and axiom reports. The final theorems depend only on
Lean's \texttt{propext}, \texttt{Classical.choice}, and \texttt{Quot.sound}.
The formal proof contains no admitted step or additional mathematical
axiom. The manuscript and its component arguments have also received
independent AI review; human mathematical review remains pending.

\begingroup
\raggedright
\bibliographystyle{plain}
\bibliography{references}
\endgroup
\end{document}
